% 第九卷 13.1.1 节：BiDAF的六层结构（正文栏宽，约108 mm）。
% 蓝色：输入与嵌入；黄色描边：编码与交互层；红色：输出分布。
% 左列为材料（长度T），右列为问题（长度J），在注意力流层汇合。
\begin{tikzpicture}[
  x=1mm,y=1mm,
  figure picture,
  emb/.style={figure node,draw=FigureInput,fill=FigureInputFill,
    minimum height=8mm,minimum width=40mm,inner ysep=2pt,
    font=\figurefont\scriptsize},
  enc/.style={figure node,draw=FigureModel,fill=FigureModelFill,
    minimum height=8mm,minimum width=40mm,inner ysep=2pt,
    font=\figurefont\scriptsize},
  wide/.style={enc,minimum width=88mm},
  res/.style={figure node,draw=FigureOutput,fill=FigureOutputFill,
    minimum height=8mm,minimum width=40mm,inner ysep=2pt,
    font=\figurefont\scriptsize},
  lab/.style={font=\figurefont\scriptsize,text=FigureMuted,anchor=east,align=right},
  fl/.style={figure flow}
]
  % 输入
  \node[font=\figurefont\scriptsize] (xin) at (22,0) {材料词$x_1,\ldots,x_T$};
  \node[font=\figurefont\scriptsize] (qin) at (70,0) {问题词$q_1,\ldots,q_J$};
  % 1-2 字符与词嵌入
  \node[emb] (xe) at (22,11) {字符CNN与词向量$\to$Highway\\$\bm X\in\mathbb R^{d\times T}$};
  \node[emb] (qe) at (70,11) {字符CNN与词向量$\to$Highway\\$\bm Q\in\mathbb R^{d\times J}$};
  % 3 上下文层
  \node[enc] (xh) at (22,25) {双向LSTM\\$\bm H\in\mathbb R^{2d\times T}$};
  \node[enc] (qu) at (70,25) {双向LSTM（共享参数）\\$\bm U\in\mathbb R^{2d\times J}$};
  % 4 注意力流层
  \node[wide,minimum height=13mm] (att) at (46,42)
    {相似矩阵$\bm S\in\mathbb R^{T\times J}$，$S_{t,j}=\bm w_{\mathrm S}^{\top}[\bm h_t;\bm u_j;\bm h_t\odot\bm u_j]$\\
     材料到问题：逐行softmax得$\widetilde{\bm U}$；\quad 问题到材料：行最大后softmax得$\widetilde{\bm H}$\\
     $\bm G_{:,t}=[\bm h_t;\tilde{\bm u}_t;\bm h_t\odot\tilde{\bm u}_t;\bm h_t\odot\tilde{\bm h}_t]$，$\bm G\in\mathbb R^{8d\times T}$};
  % 5 建模层
  \node[enc] (m) at (22,60) {两层双向LSTM\\$\bm M\in\mathbb R^{2d\times T}$};
  \node[enc] (m2) at (70,60) {再一层双向LSTM\\$\bm M_2\in\mathbb R^{2d\times T}$};
  % 6 输出层
  \node[res] (p1) at (22,75) {起点$\bm p^{(1)}=\operatorname{softmax}(\bm w_1^{\top}[\bm G;\bm M])$};
  \node[res] (p2) at (70,75) {终点$\bm p^{(2)}=\operatorname{softmax}(\bm w_2^{\top}[\bm G;\bm M_2])$};

  \draw[fl] (xin) -- (xe); \draw[fl] (qin) -- (qe);
  \draw[fl] (xe) -- (xh); \draw[fl] (qe) -- (qu);
  \draw[fl] (xh) -- (xh |- att.south); \draw[fl] (qu) -- (qu |- att.south);
  \draw[fl] (att.north -| m) -- (m);
  \draw[fl] (m) -- (m2);
  \draw[fl] (m) -- (p1);
  \draw[fl] (m2) -- (p2);
  % G 直接进入输出层
  \draw[fl,draw=FigureMuted,dash pattern=on 2pt off 1.5pt]
    (att.east) -- ++(4,0) |- (p2.east);
  \draw[fl,draw=FigureMuted,dash pattern=on 2pt off 1.5pt]
    (att.west) -- ++(-4,0) |- (p1.west);

  \node[lab] at (-1,11) {嵌入};
  \node[lab] at (-1,25) {上下文};
  \node[lab] at (-1,42) {注意力流};
  \node[lab] at (-1,60) {建模};
  \node[lab] at (-1,75) {输出};
\end{tikzpicture}
